% Shared module: interpolation, sampling, Nyquist and high frequencies.
\subsection{From scattered observations to field samples}\label{sec:alg-interp}
The input is a set of noisy observations $y_i=u(t_i,x_i)+\varepsilon_i$, $i=1,\dots,n$, with
$\varepsilon_i\sim\Normal(0,\noise^2)$, on a window $[0,T]\times[0,L)$ periodic in $x$, or a noisy field on a
regular $n_t\times n_x$ grid. Every later stage works on Fourier modes, so it needs fields on a grid.

\paragraph{Gaussian-process prior on the field}
We model $y=\bar y+f+\varepsilon$ with a constant $\bar y$ and
\begin{equation}\label{eq:gp}
  f\sim\mathcal{GP}\Big(0,\ \sum_{r=1}^{R}\eta_r^2\,k_t^{(r)}(t,t')\,k_x^{(r)}(x,x')\Big),
\end{equation}
a sum of $R$ separable Mat\'ern-$5/2$ kernels with their own amplitudes $\eta_r$ and length scales
$\ell_t^{(r)},\ell_x^{(r)}$ ($k_x$ periodic in $x$). More than one product is used because the groups we
look for have different correlation lengths: a transported profile and a diffusing one do not share a
scale. The hyperparameters $\eta=(\eta_r,\ell^{(r)},\noise,\bar y)$ are not fixed at their maximum; they are integrated out,
\begin{equation}\label{eq:marginal-eta}
  p(f\mid y)=\int p(f\mid y,\eta)\,p(\eta\mid y)\dd\eta\ \approx\ \sum_{h=1}^{N_\eta}\pi_h\,p(f\mid y,\eta_h),
  \qquad \pi_h\propto\frac{p(y,\eta_h)}{q(\eta_h)},
\end{equation}
where the $\eta_h$ are drawn from a Laplace (Gaussian) approximation $q$ of $p(\eta\mid y)$ and the importance weights
$\pi_h$ correct for the difference between $q$ and the true posterior. Each $p(f\mid y,\eta_h)$ is an ordinary Gaussian
process posterior, so the result is a weighted mixture of Gaussian processes.

\paragraph{Samples, not the mean}
The output of this stage is $\Nsamp$ fields $u^{(q)}$ on the grid, drawn from the posterior by pathwise
(Matheron) conditioning: each sample is drawn from one of the Gaussian processes $p(f\mid y,\eta_h)$
of~\eqref{eq:marginal-eta} and carries its weight, $\pi_q=\pi_h$ normalized over the samples. We do not use the posterior mean. The maps applied next
are not continuous: the interaction term is not (\cref{prop:discont}), and the factors of a numerical symbol jump
where it becomes reducible (\cref{sec:alg-structfit}), so the mean of the outputs differs from the output of the mean. In addition the mean is a Wiener filter: on a Fourier mode with
prior spectral density $S(\kappa)$ it multiplies the data by $S(\kappa)/(S(\kappa)+\noise^2/n_{\rm eff})$, with $n_{\rm eff}$ the effective number of observations per mode,
which shrinks fast-decaying, high-$\kappa$ modes towards zero. Samples keep the posterior spread;
every downstream quantity is reported as a distribution over them. On gridded input the noisy grid
itself is used, with parametric-bootstrap refits providing the spread.

\subsection{Nyquist limits and the high-frequency band}\label{sec:nyquist}
Which wavenumbers carry usable information about the symbol? Three limits apply, and the algorithm
uses all three to choose the band $\mathcal K$ of wavenumbers it analyses.

\paragraph{Spatial aliasing} On a grid of spacing $\Delta x$ the discrete Fourier transform sees
wavenumbers $\abs{\kappa}\le\kappa_x^{\max}=\pi/\Delta x$; a mode with larger $\abs\kappa$ appears at
$\kappa-2\pi q/\Delta x$ for an integer $q$. An aliased mode keeps its temporal rate $\lambda(\kappa)$
but is placed at the wrong wavenumber, so it produces a point $(\lambda,\kappa-2\pi q/\Delta x)$ that is
\emph{not} on the variety. For scattered data there is no strict Nyquist limit, but the mean spacing
$h=\sqrt{TL/n}$ plays the same role.

\paragraph{Temporal aliasing along a branch} Rates are read from a pencil of matrices built from
samples spaced by $\Delta t$, and $\log$ of an eigenvalue returns $\operatorname{Im}\lambda$ only modulo
$2\pi/\Delta t$. A branch $\tau=\Lambda_j(\xi)$ is read without ambiguity on the wavenumbers where
\begin{equation}\label{eq:nyq-t}
  \abs{\operatorname{Im}\Lambda_j(i\kappa)}<\pi/\Delta t .
\end{equation}
For transport at speed $c$ this means $\abs\kappa<\pi/(c\,\Delta t)$: the faster the transport, the
narrower the band.

\paragraph{Noise floor and decay} The observation error is modelled as Gaussian white noise with standard deviation
$\noise$ per grid point. After the normalized Fourier transform in $x$, one wavenumber bin carries noise of standard
deviation $\sigbin=\noise/\sqrt{n_x}$. A mode is counted as detected while its amplitude exceeds $z$ times this bin noise;
$z\,\sigbin$ is the \emph{detection threshold} ($z=3$ by default). A mode with initial amplitude $A_\kappa$ and rate
$\lambda(\kappa)$ has amplitude $\abs{A_\kappa}e^{\operatorname{Re}\lambda(\kappa)t}$ at time $t$, so it reaches the threshold when
\[
  \abs{A_\kappa}\,e^{\operatorname{Re}\lambda(\kappa)\,t}=z\,\sigbin,
  \qquad\text{that is}\qquad t=\frac{\log\big(\abs{A_\kappa}/(z\,\sigbin)\big)}{-\operatorname{Re}\lambda(\kappa)} .
\]
Capping by the window length and treating non-decaying modes and modes that start below the threshold gives the
time during which the mode is visible,
\begin{equation}\label{eq:visible-time}
  T_{\rm vis}(\kappa)=\min\Big\{T,\ \frac{\max\big(0,\log(\abs{A_\kappa}/(z\,\sigbin))\big)}{\max(0,-\operatorname{Re}\lambda(\kappa))}\Big\}.
\end{equation}
A rate can be estimated only if $T_{\rm vis}(\kappa)$ covers several samples; for a diffusive branch
$\operatorname{Re}\lambda=-\nu\kappa^2$, so the usable band ends near
$\kappa\approx\big(\log(\abs{A_\kappa}/z\sigbin)/(\nu\,n_{\min}\Delta t)\big)^{1/2}$, where $n_{\min}$ is the
smallest number of time samples the pencil needs. High wavenumbers of a dissipative group are
therefore lost first. This is a property of the data, not of the method; interpolation can only make
it worse, by the shrinkage described above. \Cref{fig:nyquist} shows the three limits on the running example.

\begin{figure}[tbp]
\centering
\includegraphics[width=0.66\linewidth]{fig_nyquist}
\caption{Fourier amplitudes of the worked example of \cref{sec:worked} (running example, grid, $5\%$ noise) at the first and last time, against the per-bin
detection threshold $z\,\sigbin$ ($z=3$). Only the shaded band of wavenumbers carries information
about the symbol; the spatial Nyquist limit $\pi/\Delta x$ is far away and is not the binding constraint.
What limits the band is the noise floor together with the decay of the diffusive modes
(\cref{sec:nyquist}).}\label{fig:nyquist}
\end{figure}

\paragraph{Wavenumbers per group}
The band must contain enough distinct wavenumbers on each branch for the factor to be determined.

\begin{proposition}[Counting condition]\label{prop:counting}
Let $Q=\tau-\Lambda(\xi)$ with $\Lambda$ a polynomial of degree at most $d$ with complex coefficients. The
coefficients are determined by the rates $\lambda(\kappa_\ell)=\Lambda(i\kappa_\ell)$ at $d+1$ distinct wavenumbers
$\kappa_\ell\in\R$, and not by fewer. The coefficient error is bounded by $\norm{\mathsf V^{-1}}\max_\ell\abs{\delta\lambda_\ell}$,
where $\mathsf V=\big((i\kappa_\ell)^k\big)$ is the Vandermonde matrix of the nodes; $\norm{\mathsf V^{-1}}$ grows like
$h_\kappa^{-d}$ when the nodes cluster, where $h_\kappa=\max_{\ell,\ell'}\abs{\kappa_\ell-\kappa_{\ell'}}$ is their spread. If the coefficients are real, the rates at
$-\kappa$ are the conjugates of those at $\kappa$, and $\lfloor d/2\rfloor+1$ distinct values of $\abs{\kappa}$ suffice,
at least $\lceil d/2\rceil$ of them nonzero.
\end{proposition}

\begin{proof}
$\Lambda(i\kappa)=\sum_{k\le d}\beta_k(i\kappa)^k$ is linear in $(\beta_0,\dots,\beta_d)$ with matrix $\mathsf V$,
which is invertible iff the $d+1$ nodes $i\kappa_\ell$ are distinct; with $d$ nodes a polynomial of
degree $d$ vanishing at all of them can be added. The bound is the norm of the inverse; for nodes
within $h_\kappa$ of each other the Lagrange basis polynomials have coefficients of size
$h_\kappa^{-d}$~\cite{Gau62}. For real $\beta_k$, $\operatorname{Re}\Lambda(i\kappa)=\sum_k(-1)^k\beta_{2k}\kappa^{2k}$ and
$\operatorname{Im}\Lambda(i\kappa)=\sum_k(-1)^k\beta_{2k+1}\kappa^{2k+1}$ separate the even and the odd coefficients, which are
polynomials in $\kappa^2$ with $\lfloor d/2\rfloor+1$ and $\lceil d/2\rceil$ unknowns.
\end{proof}

\roleinalg{The band $\mathcal K$ is the set of wavenumbers below both Nyquist limits whose amplitude
exceeds $z\,\sigbin$ (default $z=3$). A factor of degree $d_j$ in $\xi$ is reported as \FC{}
only if its modes occupy at least $d_j+1$ distinct signed wavenumbers $\kappa\in\mathcal K$ (for a real field the
modes at $\pm\kappa$ come in conjugate pairs, and the rule then implies the real-coefficient count of the proposition), and the whole band is
used rather than the cleanest few bins, because spreading the nodes is what keeps $\mathsf V$ well
conditioned. \Cref{fig:nyquist} shows the band for the running example.}
